\subsection{Análisis de correlación en el conjunto de datos de la flor Iris.}

\begin{enumerate}
    \item Carga el conjunto de datos de la flor Iris (disponible en MATLAB).
    \item Selecciona las cuatro variables numéricas:
    \begin{itemize}
        \item Longitud del sépalo (sepal length)
        \item Ancho del sépalo (sepal width)
        \item Longitud del pétalo (petal length)
        \item Ancho del pétalo (petal width)
    \end{itemize}
    \item Calcula la matriz de correlación entre estas variables utilizando un método adecuado (por ejemplo, el coeficiente de correlación de Pearson).
    \item Identifica pares de variables que presenten una correlación positiva significativa (valores cercanos a 1).
    \item Visualiza estas correlaciones utilizando un mapa de calor (heatmap) o gráficos de dispersión (scatter plot).
    \item Interpreta los resultados:
    \begin{itemize}
        \item ¿Qué pares de variables muestran una correlación positiva fuerte?
        \item ¿Cómo podrían estas correlaciones estar relacionadas con las características de las diferentes especies de Iris?
    \end{itemize}
    \item (Opcional) Repite el análisis separando los datos por especie (setosa, versicolor, virginica) para ver si las correlaciones varían entre ellas.
\end{enumerate}